\chapter{Testing and Validation}

\section{Testing Philosophy in Repository}
The project follows a practical mixed strategy:
\begin{itemize}[leftmargin=2em]
\item Fast unit tests for deterministic cryptographic and utility behavior.
\item Privileged integration tests gated by environment requirements.
\item Socket-level framing tests to validate protocol primitives.
\end{itemize}

This structure balances accessibility (most tests runnable without root) with realism (TUN and namespace tests where privileges permit).

\section{Crypto Test Coverage}
\texttt{tests/test\_crypto.py} covers essential security invariants:
\begin{enumerate}[leftmargin=2em]
\item Key generation length and uniqueness.
\item Base64 encode/decode round-trip.
\item Encrypt/decrypt round-trip correctness.
\item Nonce-driven ciphertext variability for same plaintext.
\item Wrong-key decryption failure.
\item Tamper detection by bit flip.
\item Invalid key length rejection.
\item Empty and large payload handling.
\end{enumerate}

These are strong educational tests because they validate both nominal behavior and negative cases that represent real-world attack surfaces.

\section{Tunnel Tests}
\texttt{tests/test\_tunnel.py} includes privilege-aware checks:
\begin{itemize}[leftmargin=2em]
\item TUN creation and context manager semantics.
\item Timeout behavior for non-blocking reads.
\item Pre-open error path assertions.
\item Convenience helper coverage.
\end{itemize}

Root and \texttt{/dev/net/tun} checks are used to skip unsupported environments cleanly.

\section{End-to-End and Integration Tests}
\texttt{tests/test\_end2end.py} validates:
\begin{itemize}[leftmargin=2em]
\item Config template generation correctness.
\item Shared key consistency between endpoint configs.
\item Basic TCP socket message exchange.
\item Framed packet round-trip using protocol helpers.
\item Optional namespace setup script behavior (root-gated).
\end{itemize}

\section{Validation Matrix}
\begin{table}[H]
\centering
\caption{Validation Matrix by Subsystem}
\begin{tabular}{p{0.30\textwidth}p{0.18\textwidth}p{0.42\textwidth}}
\toprule
\textbf{Subsystem} & \textbf{Coverage Level} & \textbf{Notes} \\
\midrule
Crypto primitives & High & Includes tamper, wrong key, nonce uniqueness patterns \\
Protocol framing & Medium-High & Socket loopback framing tests included \\
TUN lifecycle & Medium & Privileged tests with environment guards \\
Namespace scripts & Medium & Some script invocation checks under root \\
Full tunnel path & Medium & Demonstrated via scripts and manual ping workflows \\
\bottomrule
\end{tabular}
\end{table}

\section{Testing Gaps}
Despite good baseline coverage, several gaps remain:
\begin{itemize}[leftmargin=2em]
\item No sustained soak tests for long-running stability.
\item No throughput or latency benchmark suite.
\item No fuzz testing for framing parser robustness.
\item No adversarial replay simulation.
\item No CI matrix across Linux distributions/kernels.
\end{itemize}

\section{Suggested Test Roadmap}
\begin{enumerate}[leftmargin=2em]
\item Add pytest markers for \texttt{unit}, \texttt{integration}, \texttt{privileged}, \texttt{network}.
\item Add synthetic packet replay tests once packet IDs are introduced.
\item Introduce stress harness for packet flood and disconnect/reconnect scenarios.
\item Capture benchmark outputs (pps, Mbps, latency percentile).
\item Run shell lint and static checks for script quality.
\end{enumerate}

\section{Reproducibility Notes}
To make results reproducible in reports or classrooms, record:
\begin{itemize}[leftmargin=2em]
\item Kernel version and distro.
\item Python version and dependency lock state.
\item Privilege model (root/container/sudo policy).
\item Namespace setup mode used.
\end{itemize}

\section{Trust Boundaries of Tests}
Passing tests should be interpreted correctly:
\begin{itemize}[leftmargin=2em]
\item They validate implementation consistency, not production threat resistance.
\item They increase confidence in pedagogical correctness.
\item They do not replace formal cryptographic protocol proofs.
\end{itemize}

\section{Conclusion}
Testing quality is strong relative to project size and educational intent. By adding performance and adversarial test dimensions, WhisperTunnel can evolve from a teaching artifact to a robust experimental platform.
